% Generated draft: requires mathematical and source review.
\documentclass[11pt]{article}
\usepackage[margin=25mm]{geometry}
\usepackage{fontspec}
\setmainfont{Latin Modern Roman}
\usepackage{amsmath,amssymb,mathtools,unicode-math}
\setmathfont{Latin Modern Math}
\usepackage{enumitem,needspace,xurl,hyperref}
\setlength{\parindent}{0pt}
\setlength{\parskip}{5pt}
\setlength{\emergencystretch}{3em}
\newcommand{\N}{\mathbb{N}}
\newcommand{\Z}{\mathbb{Z}}
\newcommand{\Q}{\mathbb{Q}}
\newcommand{\R}{\mathbb{R}}
\newcommand{\C}{\mathbb{C}}
\newcommand{\F}{\mathbb{F}}
\newcommand{\A}{\mathbb{A}}
\newcommand{\PP}{\mathbb{P}}
\newcommand{\E}{\mathbb{E}}
\newcommand{\eps}{\varepsilon}
\newcommand{\dd}{\,\mathrm{d}}
\begin{document}
\section*{UnsolvedMath ID: 30003532}
\textbf{Source problem code:} catalogue entry
\subsection{Definitions and mathematical statement}
Let \(M=\mathbb{CP}^{2}\), the space of complex one-dimensional linear subspaces of \(\C^{3}\), and let \(\omega\) denote the intended symplectic form. The supplied record does not specify this form; the conventional choice is the Fubini--Study form, with normalization still to be specified. For a fixed \(\omega\), a smooth Hamiltonian diffeomorphism is a map \(\phi=\phi_H^1\), where \(\phi_H^t\) is the flow of a smooth time-dependent Hamiltonian \(H\colon[0,1]\times M\to\R\), with vector field determined by \(\iota_{X_{H_t}}\omega=-\dd H_t\). Write \(\operatorname{Ham}(M,\omega)\) for the set of these diffeomorphisms and define
\[
\operatorname{Per}(\phi)=\bigcup_{k\geq1}\operatorname{Fix}(\phi^k),\qquad
\operatorname{Fix}(\phi^k)=\{x\in M:\phi^k(x)=x\}.
\]
Thus ``exactly three periodic points'' means
\[
\#\operatorname{Per}(\phi)=3,
\]
counting points over all positive integer periods, rather than counting periodic orbits or only fixed points. The source calls the desired maps Hamiltonian pseudo-rotations; it does not specify additional nondegeneracy conditions.

The construction question is whether the Anosov--Katok conjugation method can produce a map
\[
\phi\in\operatorname{Ham}(\mathbb{CP}^{2},\omega),\qquad \#\operatorname{Per}(\phi)=3,
\]
that is a genuine pseudo-rotation beyond true rotations. For reference, with the Fubini--Study form, the usual projective unitary rotations have the form
\[
R_U([z])=[Uz],\qquad U\in\operatorname{U}(3).
\]
The record does not specify whether ``beyond true rotations'' excludes only maps of this form, or also maps conjugate to them, nor does it specify the permitted conjugacies. Consequently, the displayed periodic-point condition is precise, but the required exclusion of rotations and the exact admissible Anosov--Katok construction remain unresolved parts of the formulation. The smooth setting above is the usual Hamiltonian interpretation; the source does not explicitly state the regularity or convergence topology required of the construction.
\subsection{Short English statement}
Can the Anosov--Katok conjugation method construct a Hamiltonian diffeomorphism of \(\mathbb{CP}^{2}\) having exactly three periodic points in total and satisfying the source's requirement of being beyond true rotations? The symplectic convention, construction regularity, and meaning of exclusion of rotations require clarification.
\subsection{Sources}
\begin{itemize}
\item Local UnsolvedMath registry, supplied complete source record: permanent ID \(30003532\), exact source problem code \texttt{catalogue entry}, title ``Hamiltonian Pseudo-Rotations of the Complex Projective Plane.'' The supplied proposed year is 2017; no individual proposer is identified.
\item Supplied report citation: \emph{Dynamische Systeme (2018)}, Oberwolfach Reports, pp.\ 1987--2046, DOI \texttt{10.4171/owr/2017/32}. \url{https://doi.org/10.4171/owr/2017/32}. The literature-source list labels this same reference ``Oberwolfach Report 32/2017.'' The differing year descriptions are preserved without reconciliation.
\item Ginzburg--Gurel, \emph{Hamiltonian pseudo-rotations}. Supplied literature link: \url{https://arxiv.org/abs/1712.09766}.
\item Ginzburg--Gurel, \emph{Pseudo-rotations versus rotations}. Supplied literature link: \url{https://arxiv.org/abs/1812.05782}.
\item Literature assessment dated exactly \texttt{2026-08-21}: the supplied curation labels the problem open and states that Floer-theoretic results establish invariant sets and recurrence properties for Hamiltonian pseudo-rotations of complex projective spaces, while the requested general construction and the existence of ghost pseudo-rotations remain open. This is supplied dated curation derived from the OWR material, not an independently verified assessment. The record explicitly requests independent verification of time-sensitive status and cited claims.
\item The supplied statement-status label is \texttt{corrected\_verified}, accompanied by a claim of manual semantic, typographic, and LaTeX review. This records the supplied editorial assessment and does not establish mathematical correctness or independent source verification.
\end{itemize}
\end{document}
